\chapter{Step 2: images into a cache}
\label{ch:cache}
\whereami{3}

\begin{tldr}
\begin{itemize}
\item Reading DICOM files is slow, so we did it \textbf{once} and saved every study as three small 3D arrays: the
      \textbf{cache}.
\item Per study: pick 3 series (sagittal, coronal, axial, fluid-sensitive), \textbf{sort} the slices,
      \textbf{rotate/flip} them to one convention, \textbf{normalise} brightness, \textbf{crop} 140\,mm around the knee,
      \textbf{resize} to 24 slices of 256\,$\times$\,256 pixels.
\item The same Python file prepares the hidden test images inside the submission, so training and test images are
      made \textbf{identically}.
\end{itemize}
\end{tldr}

\section{Why a cache?}

Training reads every study 10 times per model, and we trained 5 models per experiment. Decoding hundreds of thousands of
DICOM files each time would take hours. Decoding once (\ca 1.5--2 hours of CPU) and saving compact arrays makes every
later experiment fast.

\begin{figure}[H]
\centering
\begin{tikzpicture}[x=0.95mm,y=1mm,
  st/.style={draw=accent, fill=keyfill, rounded corners=1.5mm, minimum width=20mm, minimum height=13mm,
             align=center, font=\scriptsize, text width=19mm},
  ar/.style={-{Stealth[length=1.6mm]}, line width=0.7pt, accent}]
  \foreach \i/\t in {0/{\textbf{a} pick\\ 3 series}, 1/{\textbf{b} sort\\ slices}, 2/{\textbf{c} same\\ orientation},
                     3/{\textbf{d} note\\ left/right}, 4/{\textbf{e} normalise\\ brightness}, 5/{\textbf{f} crop\\ the knee},
                     6/{\textbf{g} resize\\ 24$\times$256$\times$256}} {
    \node[st] (n\i) at ({\i*22.5},0) {\t};
  }
  \foreach \i [evaluate=\i as \j using int(\i+1)] in {0,...,5} \draw[ar] (n\i) -- (n\j);
  \node[font=\scriptsize, text=muted] at (0,-11) {\faFolder\ DICOM files};
  \node[font=\scriptsize, text=muted] at (135,-11) {\faCube\ 3 arrays per study};
\end{tikzpicture}
\caption{The seven steps, applied to every study. The code lives in one file, \code{knee\_preproc.py}, saved together
with the cache.}
\label{fig:cacheflow}
\end{figure}

\section{a. Pick one series per role}

A study has about 5.5 series. We keep three \term{roles}: \code{sag\_fs}, \code{cor\_fs}, \code{ax\_fs} (sagittal,
coronal, axial; fluid-sensitive). \code{train\_series.csv} says what each series is; a small score picks the best:

\begin{figure}[H]
\centering
\begin{tikzpicture}[x=1mm,y=1mm,
  ser/.style={draw=grid, fill=white, rounded corners=1mm, minimum width=34mm, minimum height=8mm, font=\scriptsize,
              align=left, anchor=west},
  win/.style={ser, draw=us, line width=0.9pt, fill=keyfill}]
  \node[ser] at (0,20) {Sagittal, fluid, fat-sat \hfill \textbf{7}};
  \node[ser] at (0,10) {Sagittal, T1 \hfill 1};
  \node[ser] at (0,0)  {Coronal, fluid, fat-sat \hfill \textbf{7}};
  \node[ser] at (0,-10) {Coronal, T1 \hfill 1};
  \node[ser] at (0,-20) {Axial, fluid, fat-sat \hfill \textbf{7}};
  \node[win] (r1) at (70,20) {\code{sag\_fs}};
  \node[win] (r2) at (70,0) {\code{cor\_fs}};
  \node[win] (r3) at (70,-20) {\code{ax\_fs}};
  \foreach \y in {20,0,-20} \draw[-{Stealth}, us, line width=0.8pt] (35,\y) -- (69,\y);
  \node[font=\scriptsize, align=left, anchor=west] at (108,0) {score =\\ \hspace*{2mm}4 if fluid-sensitive\\
     \hspace*{2mm}+ 2 if also fat-suppressed\\ \hspace*{2mm}+ 1 if 12--80 slices\\[1mm] highest score wins;\\
     no fluid series in a plane:\\ take any (``fallback'')};
\end{tikzpicture}
\caption{Series selection for one study. All three roles were available in 100\,\% of training studies.}
\end{figure}

\section{b. Sort the slices}

File names are random IDs. The only reliable order is the physical one: project each slice's
\code{ImagePositionPatient} onto the slice direction and sort by that number.

\begin{figure}[H]
\centering
\begin{tikzpicture}[x=1mm,y=1mm]
  \foreach \k/\o in {0/3,1/0,2/5,3/1,4/4,5/2} {
    \pgfmathsetmacro{\g}{35+\o*9}
    \fill[dark!\g] ({\k*9},0) rectangle ++(7,10);
    \node[font=\tiny] at ({\k*9+3.5},-2.5) {\o};
  }
  \node[font=\scriptsize, align=center] at (25,-8) {sorted by file name: jumbled};
  \draw[-{Stealth}, line width=1pt, muted] (58,5) -- (72,5);
  \foreach \k in {0,...,5} {
    \pgfmathsetmacro{\g}{35+\k*9}
    \fill[dark!\g] ({78+\k*9},0) rectangle ++(7,10);
    \node[font=\tiny] at ({78+\k*9+3.5},-2.5) {\k};
  }
  \node[font=\scriptsize, align=center] at (103,-8) {sorted by position (mm): correct};
\end{tikzpicture}
\caption{Wrong slice order turns a 3D stack into noise for anything that looks at neighbouring slices.}
\end{figure}

\section{c. One orientation for everybody}

Sites store the same plane in different ways: rotated, mirrored, upside down. The header's
\code{ImageOrientationPatient} says how, so every series can be flipped or transposed into one convention (for
sagittal: front of the knee on the left, head at the top).

\begin{figure}[H]
\centering
\begin{tikzpicture}
  \foreach \x/\lab in {0/{site A}, 2.3/{site B}, 4.6/{site C}, 6.9/{site D}} {
    \draw[grid] (\x,0) rectangle ++(2,2.3);
    \node[font=\scriptsize] at (\x+1,-0.3) {\lab};
  }
  \begin{scope}[shift={(0.05,0.05)}, scale=0.19] \kneesagittal \end{scope}
  \begin{scope}[shift={(4.25,0.05)}, xscale=-0.19, yscale=0.19] \kneesagittal \end{scope}
  \begin{scope}[shift={(4.65,2.25)}, xscale=0.19, yscale=-0.19] \kneesagittal \end{scope}
  \begin{scope}[shift={(8.9,0.28)}, rotate=90, scale=0.165] \kneesagittal \end{scope}
  \draw[-{Stealth}, line width=1pt, muted] (9.3,1.15) -- (10.6,1.15);
  \draw[accent, line width=1pt] (11,0) rectangle ++(2,2.3);
  \begin{scope}[shift={(11.05,0.05)}, scale=0.19] \kneesagittal \end{scope}
  \node[font=\scriptsize\bfseries, text=accent] at (12,-0.3) {canonical};
\end{tikzpicture}
\caption{Four ways the same sagittal image can be stored, and the one convention every series is turned into. We
checked that five real storage variants of one series produced identical output.}
\end{figure}

\section{d. Left or right knee: written down, not changed}

For a right knee the medial side is towards the patient's left, for a left knee the other way. The cache keeps the
patient's own orientation and records the \code{Laterality} tag (present in 49.8\,\% of studies) so a later step can
mirror left knees if wanted. We tried that in Chapter~\ref{ch:experiments} (``v7'').

\section{e. Normalise the brightness}

\begin{figure}[H]
\centering
\begin{tikzpicture}
\begin{axis}[width=10cm, height=4.6cm, xmin=0, xmax=3000, ymin=0, ymax=1.15, ytick=\empty, axis y line=none,
    axis x line*=bottom, xlabel={raw pixel value of one series}, xlabel style={font=\footnotesize},
    xticklabel style={font=\scriptsize}, scaled x ticks=false]
  \addplot[fill=us!25, draw=us, domain=0:3000, samples=150] {1.0*exp(-((x-40)/45)^2) + 0.45*exp(-((x-700)/220)^2) + 0.25*exp(-((x-1400)/260)^2) + 0.03*exp(-((x-2600)/200)^2)} \closedcycle;
  \draw[them, line width=1.2pt, dashed] (axis cs:15,0) -- (axis cs:15,1.1) node[above, font=\scriptsize] {0.5\,\%};
  \draw[them, line width=1.2pt, dashed] (axis cs:2150,0) -- (axis cs:2150,1.1) node[above, font=\scriptsize] {99.5\,\%};
  \node[font=\scriptsize, align=left, anchor=west] at (axis cs:120,0.8) {air (background)};
  \node[font=\scriptsize, align=center] at (axis cs:1050,0.6) {tissue};
  \node[font=\scriptsize, align=center] at (axis cs:2600,0.3) {rare bright\\ outliers};
\end{axis}
\end{tikzpicture}
\caption{Per series: find the 0.5th and 99.5th percentile, map that range to 0--255, clip the rest. Outliers no longer
decide the contrast, and every series ends up on the same scale (compare Chapter~\ref{ch:mri}).}
\end{figure}

\section{f. Crop the knee}

The field of view varies between sites (often 150--200\,mm), and much of it is thigh and shin. We found the knee
automatically and cut a fixed square around it:

\begin{figure}[H]
\centering
\begin{tikzpicture}[scale=0.5]
  \fill[black] (-1.5,-1.0) rectangle (11.5,13.0);
  \begin{scope}[shift={(0,0)}]
    \clip (-1.5,-1.0) rectangle (11.5,13.0);
    \kneecoronal
  \end{scope}
  \draw[us, line width=1.6pt, dashed] (0.5,0.5) rectangle (9.5,9.5);
  \node[font=\scriptsize\bfseries, text=us, fill=white, inner sep=1pt, anchor=south west] at (0.5,9.6) {140\,mm (v5)};
  \draw[them, line width=1.6pt] (2.1,2.3) rectangle (7.9,8.1);
  \node[font=\scriptsize\bfseries, text=them, fill=white, inner sep=1pt, anchor=north west] at (2.1,2.2) {90\,mm (v8)};
  \node[font=\scriptsize\bfseries, text=white, anchor=north west] at (-1.4,12.9) {full field of view, e.g.\ 200\,mm};
\end{tikzpicture}\hspace{6mm}
\begin{minipage}[b]{0.42\linewidth}\small
\textbf{How the knee is found:}
\begin{enumerate}
\item average all slices of the series into one image;
\item ``tissue'' = pixels brighter than 30\,\% of the 99th percentile;
\item the tissue box = rows and columns that are at least 5\,\% tissue;
\item centre a 140\,mm square on that box (pad with black if the image is smaller).
\end{enumerate}
We stored how much tissue fell outside (\code{tissue\_outside}) to check that no knee was cut.\par\medskip
\end{minipage}
\caption{The crop window. The bigger the window, the less detail per millimetre after resizing to 256 pixels.}
\label{fig:crop}
\end{figure}

\section{g. Resize to one shape}

Every series becomes exactly \textbf{24 slices of 256\,$\times$\,256 pixels}, stored as 8-bit integers (0--255):
linear interpolation between slices, smooth (anti-aliased) resizing within a slice.

\begin{center}\footnotesize
\begin{tabular}{lllll}
\toprule
Cache & Window & Pixels & mm/pixel & Result (Chapter~\ref{ch:experiments}) \\
\midrule
first (v1--v4) & whole field of view & 224 & \ca 0.7--0.9 & 0.905 public \\
\textbf{v5} & 140\,mm around the knee & 256 & 0.55 & \textbf{0.918} public \\
v6 & 130\,mm & 320 & 0.41 & +0.002--0.004 val.: stopped \\
v8 & 90\,mm on the joint & 256 & 0.35 & prepared; result not recorded \\
\midrule
\textcolor{them}{winners} & 140\,mm; 112\,mm on the joint & 384--512 & 0.27--0.36 & 0.960 private \\
\bottomrule
\end{tabular}
\end{center}

\begin{innumbers}
One study = 3 roles $\times$ 24 $\times$ 256 $\times$ 256 bytes $\approx$ 4.7\,MB, so the whole training set is
\ca 21\,GB in memory. Kaggle notebooks may output at most 500 files and 20\,GB, so the cache was written as
\textbf{45 compressed shards} of 100 studies each.
\end{innumbers}

\section{The same code at test time}

\begin{keyidea}
The submission notebook imports the \emph{same} \code{knee\_preproc.py} from the cache dataset and runs it on the
hidden test DICOMs. Any difference between training and test preparation (another resize method, another
orientation rule) would silently cost score, and you would never see why.
\end{keyidea}

\begin{watchout}
24 slices per series is a compromise. Native series have 20--40 slices; resampling to 24 blurs thin structures
across slices. The winners kept up to 64 \emph{native} slices and higher in-plane resolution, and used five series
instead of three. That costs much more GPU time, which on Kaggle's free T4s we did not have.
\end{watchout}

\begin{checkpoint}
\begin{itemize}
\done I can list the seven preparation steps and say why each is needed.
\done I know why the cache module is reused in the submission.
\quiz Q8.1: Why normalise each series separately instead of the whole dataset at once?
\quiz Q8.2: Cropping from 200\,mm to 140\,mm and keeping 256 pixels: how much finer does each pixel get?
\quiz Q8.3: The cache does not mirror left knees. What would go wrong if it mirrored only the sagittal series?
\end{itemize}
\end{checkpoint}
